% Track e1 -- One particle or two: the second rung Cl(6,6).
% Verification: investigations/emergence/verify/e1_two_from_one.py (claims e1-1a ... e1-5f).
% New references: sections/bib_e1.tex.
\section{One particle or two: the second rung}
\label{e1:sec:main}

The conventions of \cite{ES17} say that a constituent is a tensor factor $\Cl(3,3)$ of the tower, that the word never refers to a
particle, and that three constituents do not describe three particles (Section~1.4 of \cite{ES17}). The duplication
$\Cl(3,3)\to\Cl(6,6)$ of \cite{RE} (Sections~7--9 there) is therefore not yet the emergence of two particles from one. This section
asks in what precise sense, and under which added structures, the second rung $\Cl(6,6)$ carries two particles rather than one
particle with a larger internal space. Four tests separate the two readings: whether the copies have separate mass shells, whether
there is one complex structure under which both copies are complex Dirac spinors, whether some operator counts particles, and whether
exchange of the copies sorts states by statistics. The answer has three parts. The graded tower describes one particle: its two mass
terms anticommute, they add in quadrature, and no state has a sharp rest energy in each copy (fixed). A two-particle reading exists, and
it is fixed as an algebra: the commutant of the first copy is a second, commuting $\Cl(3,3)$, whose generators are the graded ones
dressed by the chirality of the first copy (fixed). But it needs three further structures, and none of them is selected. They are
a common complex structure, which exists only on a $32$-dimensional subspace, an additive Hamiltonian that is not a vector of
$\Cl(6,6)$, and an ungraded exchange that is not the Clifford exchange (constructed; the selection of a complex structure is obstructed
under the diagonal action). The Fock number of the tower counts quanta, not particles. No gauge-admissible field that commutes with the
binding root changes it (fixed). It is not invariant under the interior Lorentz algebra, under the boosts of $\lG$, or under the
charge $J$ (fixed), so as a particle number it is an internal reading (interpretive). Identifiers e1-$n$ refer to the \texttt{PASS} lines of
\texttt{verify/e1\_two\_from\_one.py} (\cref{e1:tab:claims}).

\paragraph{Notation.} $S=\R^8$ with the realisation of equation~(1) of \cite{ES17}, $\omega=e_1\cdots e_6$, $J=e_2e_3$, $T=e_1e_2e_3$,
$\Omega=e_1e_4e_5e_6$, and $\gamma^\mu=(e_1,e_4,e_5,e_6)$. Appendix~A of \cite{ES17} records $\omega=J\Omega=\Gamma$, the chirality. On
$S_2=S\otimes S=\R^{64}$ the graded generators are $E_{(1,a)}=e_a\otimes\id$ and $E_{(2,a)}=\omega\otimes e_a$ (equation~(2) of \cite{ES17}),
and the \emph{ungraded} copy-2 generators are $u_a=\id\otimes e_a$. We write $\Gamma^{(1)}=\omega\otimes\id$ and $\Gamma^{(2)}=\id\otimes\omega$
for the volume elements of the two constituents, $J^{(1)}=J\otimes\id$, $J^{(2)}=\id\otimes J$, $T^{(1)}=T\otimes\id$, $T^{(2)}=\id\otimes T$,
and $\omega_S=E_{(1,2)}E_{(1,3)}E_{(2,2)}E_{(2,3)}$ for the volume element of the four surplus directions. On vectors, $K_n$ is the
para-complex structure $Ke_i=e_{i+3}$, $Ke_{i+3}=e_i$ on each constituent; it is $-\cK$ of Appendix~A of \cite{ES17}. $G_n$ is the binding
root, and $\lG$ and $\lint$ are the field-strength and interior Lorentz algebras of one constituent. The Fock modes are
$a^\pm_{k,i}=\half(E_{(k,i)}\pm E_{(k,i+3)})$, with $N=\sum_{k,i}a^+_{k,i}a^-_{k,i}$ and $N^{(k)}$ the part of constituent $k$
(Proposition~7.2 of \cite{RE}). The rest masses are $m_1,m_2\ge0$ and $M=(m_1^2+m_2^2)^{1/2}$.

\subsection{Two copies of \texorpdfstring{$\Cl(3,3)$}{Cl(3,3)} inside \texorpdfstring{$\Cl(6,6)$}{Cl(6,6)}}
\label{e1:sec:copies}

\begin{proposition}[The commuting copy]
\label{e1:prop:copies}
\begin{enumerate}[label=(\roman*),leftmargin=2.2em]
\item The twelve $E_{(k,a)}$ satisfy the Clifford relations of $\R^{6,6}$. The six $u_a$ satisfy those of $\R^{3,3}$ and commute with every
$E_{(1,b)}$. Moreover $u_a=\Gamma^{(1)}E_{(2,a)}=E_{(1,1)}\cdots E_{(1,6)}E_{(2,a)}$, a blade of grade $7$.
\item The commutant of $\Cl(3,3)\otimes\id$ in $M_{64}(\R)$ is spanned by $64$ blades of $\Cl(6,6)$, each equal to $\pm\id\otimes e_J$; it is
$\id\otimes M_8(\R)$, the algebra generated by the $u_a$. A second copy of $\Cl(3,3)$ commuting with the first is therefore unique up to
conjugation by $\id\otimes GL(8,\R)$.
\item The dressing is the chirality and the Fock parity of the first copy: $\Gamma^{(1)}=(J\Omega)\otimes\id=-(-1)^{N^{(1)}}$.
\item Both the ungraded product $\Cl(3,3)\otimes\Cl(3,3)$ and the graded product $\Cl(3,3)\hat\otimes\Cl(3,3)=\Cl(6,6)$ are $M_{64}(\R)$, with
the same $4096$ blades as basis. The twelve elements $e_a\otimes\id$, $u_a$ span no quadratic space: $(e_1\otimes\id+u_1)^2=2\,\id+2\,e_1\otimes e_1$.
\item At rung $n$ the ungraded generator of constituent $k$, $\id^{\otimes(k-1)}\otimes e_a\otimes\id^{\otimes(n-k)}$, equals
$\Gamma^{(1)}\cdots\Gamma^{(k-1)}E_{(k,a)}$, a blade of grade $6(k-1)+1$. Ungraded generators of distinct constituents commute and graded
ones anticommute (checked at $n=3,4$).
\end{enumerate}
\end{proposition}

\begin{proof}
(i) $(\omega\otimes\id)(\omega\otimes e_a)=\omega^2\otimes e_a=\id\otimes e_a$, and $\Gamma^{(1)}=E_{(1,1)}\cdots E_{(1,6)}$. (ii) Conjugation by an
invertible blade multiplies every blade by $\pm1$ (equation~(3) of \cite{ES17}), so the commutant of a set of blades is spanned by the
blades it contains. A blade commutes with all $E_{(1,b)}$ exactly when it is $\Gamma^{(1)\varepsilon}$ times a blade of copy~2 with grade
parity $\varepsilon$, that is $\pm\id\otimes e_J$ (Proposition~7.1(iii) of \cite{RE}). The uniqueness is the Skolem--Noether theorem in
$\id\otimes M_8(\R)\cong M_8(\R)$. (iii) $\omega=J\Omega$ by moving $e_1$ past $e_2e_3$, and $(-1)^{N_1+N_2+N_3}=-\omega$ is Proposition~7.2
of \cite{RE}. (iv) Distinct blades $B,B'$ have $\Tr(B^{\!\top}B')=0$ because $B^{\!\top}B'$ is $\pm$ a blade other than $\id$. The square is
a one-line expansion. (v) $\Gamma^{(j)}=\id^{\otimes(j-1)}\otimes\omega\otimes\id^{\otimes(n-j)}$, and the $\omega$'s cancel in pairs.
\end{proof}

Parts (i)--(iii) restate Proposition~7.1 of \cite{RE} and Appendix~A of \cite{ES17} (fixed there), with the ungraded copy named. The
Klein--Jordan--Wigner reading is inherited \cite{JordanWigner28,Klein38}, and so is the exponential law \cite{Chevalley54,LawsonMichelsohn89}. The
identification of the Klein factor with the chirality of the first copy, $u_a=\Gamma^{(1)}E_{(2,a)}$, is fixed (e1-1a to e1-1d). Two
remarks organise the rest of the section. First, bivectors do not see the grading: $E_{(2,a)}E_{(2,b)}=\id\otimes e_ae_b$. Hence each
constituent's $\lint^{(k)}$ and $\lG^{(k)}$, its charge $J^{(k)}=E_{(k,2)}E_{(k,3)}$ and its number $N^{(k)}$ are the same operators in both
readings, and the two Lorentz algebras commute in both. Second, the readings differ only in their odd elements, the gamma matrices and
the mass terms. The algebra $M_{64}(\R)$ is the same; what changes is which of its elements are taken as the vectors.

\subsection{Rest energies: quadrature against addition}
\label{e1:sec:rest}

Put $H_g=m_1E_{(1,1)}+m_2E_{(2,1)}$, the rest Hamiltonian built from the tower's vectors, and $H_u=m_1\,e_1\otimes\id+m_2\,\id\otimes e_1$,
built from the commuting copy. On one constituent the rest flow of the real Dirac form is $\dot\Psi=-JH(0)\Psi$ with $JH(0)=mT$, of angular
frequency $m$ (Section~4.4 of \cite{ES17}; Proposition~4.9 of \cite{RE}).

\begin{proposition}[Graded and ungraded rest energies]
\label{e1:prop:rest}
\begin{enumerate}[label=(\roman*),leftmargin=2.2em]
\item $H_g^2=M^2\,\id$ and $\Tr H_g=0$, so $\spec H_g=\{\pm M\}$, each of multiplicity $32$.
\item With $B=E_{(2,1)}E_{(1,1)}$, $B^2=-\id$, and $U_\phi=\cos\phi+\sin\phi\,B$, $\tan2\phi=m_2/m_1$: $U_\phi(ME_{(1,1)})U_\phi^{-1}=H_g$. The graded
rest Hamiltonian is one Dirac mass along a unit vector of $\R^{6,6}$, turned by a compact element of $\mathrm{Spin}(6,6)$.
\item $\spec H_u=\{\pm m_1\pm m_2\}$, each of multiplicity $16$. $H_u=m_1E_{(1,1)}+m_2\,\Gamma^{(1)}E_{(2,1)}$ has components of grade $1$ and $7$.
For $m_1\neq m_2$, both nonzero, it has four distinct eigenvalues, so it is not conjugate to any vector of $\R^{6,6}$, which satisfies
$v^2=\eta_2(v,v)\id$ and has at most two.
\item Joint mass shells. No nonzero $\psi$ satisfies $E_{(1,1)}\psi=s_1\psi$ and $E_{(2,1)}\psi=s_2\psi$ with $s_1,s_2=\pm1$. Each joint sign
$(s_1,s_2)$ of $e_1\otimes\id$ and $\id\otimes e_1$ occurs on $16$ dimensions. More generally, the graded Dirac operators
$D_k=\sum_\mu p^{(k)}_\mu E_{(k,\mu)}$ anticommute, with $(D_1+D_2)^2=(p^{(1)}\cdot p^{(1)}+p^{(2)}\cdot p^{(2)})\id$, while their ungraded
counterparts commute.
\item Rest flows. $J^{(1)}$ and $J^{(2)}$ commute with $H_g$ and $H_u$. The ungraded generator $Z_u=m_1T^{(1)}+m_2T^{(2)}$ is a sum of two commuting
complex structures, with eigenvalues $i(\pm m_1\pm m_2)$: $e^{-\lambda Z_u}$ winds a two-torus with angular frequencies $m_1$ and $m_2$. The
graded generator $Z_g=J^{(1)}H_g$ has $Z_g^2=-M^2\id$, so $e^{-\lambda Z_g}=\cos(M\lambda)-\sin(M\lambda)Z_g/M$ has the single frequency $M$.
\item At rung $n$ the graded $E_{(k,1)}$ anticommute and the ungraded ones commute, with traceless products over every nonempty subset. Hence
$\spec\bigl(\sum_km_kE_{(k,1)}\bigr)=\{\pm(\sum_km_k^2)^{1/2}\}$, while the ungraded sum has spectrum $\{\sum_k\pm m_k\}$, each value of
multiplicity $4^n$.
\end{enumerate}
\end{proposition}

\begin{proof}
(i) $E_{(1,1)}$ and $E_{(2,1)}$ are anticommuting traceless involutions. (ii) $U_\phi XU_\phi^{-1}=\cos^2\!\phi\,X+\sin\phi\cos\phi\,[B,X]-\sin^2\!\phi\,BXB$, with
$[B,E_{(1,1)}]=2E_{(2,1)}$ and $BE_{(1,1)}B=E_{(1,1)}$. (iii) Commuting involutions with traceless products have all four joint signs with equal
multiplicity. (iv) If $E_{(1,1)}\psi=s_1\psi$ and $E_{(2,1)}\psi=s_2\psi$, then $s_1s_2\psi=E_{(1,1)}E_{(2,1)}\psi=-E_{(2,1)}E_{(1,1)}\psi=-s_1s_2\psi$. Generators of
distinct constituents anticommute, so the cross terms of $(D_1+D_2)^2$ cancel. (v) $T^{(k)}=J^{(k)}\times(\text{the mass term of copy }k)$, both
$T^{(k)}$ square to $-\id$, and they commute. $J^{(1)}$ commutes with $E_{(1,1)}$ and $E_{(2,1)}$ and squares to $-\id$, so
$(J^{(1)}H_g)^2=-H_g^2$. (vi) As (i) and (iii), with $2^n$ joint signs.
\end{proof}

Parts (i) and (iii) at $m_1=m_2$ restate Proposition~8.11(iv) of \cite{RE}. The general masses, the rotation (ii), the joint mass shells
(iv), the flows (v) and the ladder (vi) are fixed (e1-2a to e1-2g). The free-particle facts behind (iv) are inherited. A system of two
particles has two mass-shell conditions, one per particle, imposed simultaneously, as in the two-body Dirac (Breit) equation
\cite{Breit29}. A single particle in a doubled momentum space has one, $p^{(1)}\cdot p^{(1)}+p^{(2)}\cdot p^{(2)}=M^2$. By (ii) and (iv)
the graded tower is of the second kind. Its two copies have no joint mass shell, its rest energy is one value $M$, and the second copy
enters the one particle as further momentum and mass directions. The ungraded operators are of the first kind: the energies add, and
the pair threshold is $m_1+m_2$.

\subsection{Complex structures: the two-particle space and the selected unit}
\label{e1:sec:complex}

The paper reads $J=e_2e_3$ as the imaginary unit of one constituent (Theorem~4.1 of \cite{ES17}). A complex two-particle space needs one
imaginary unit acting as $J$ on both factors.

\begin{proposition}[The complex two-particle space]
\label{e1:prop:W}
\begin{enumerate}[label=(\roman*),leftmargin=2.2em]
\item $J^{(1)}$ and $J^{(2)}=E_{(2,2)}E_{(2,3)}$ commute and square to $-\id$, and $J^{(1)}J^{(2)}=\omega_S$ is a symmetric traceless involution. The
subspace $W=\{\psi:J^{(1)}\psi=J^{(2)}\psi\}=\ker(\omega_S+\id)$ has dimension $32$, and its orthogonal complement is
$\overline W=\ker(J^{(1)}+J^{(2)})=\ker(\omega_S-\id)=\operatorname{im}(J^{(1)}-J^{(2)})$. With $J:=J^{(1)}|_W=J^{(2)}|_W$ as imaginary unit,
$W\cong S\otimes_\C S$, and $\overline W\cong S\otimes_\C\overline S$.
\item A blade of $\Cl(6,6)$ with an even number of surplus indices preserves $W$ and $\overline W$. A blade with an odd number exchanges them,
for instance $\id\otimes e_2$ (charge conjugation $C=e_2$ of copy~2 alone, Theorem~5.1 of \cite{ES17}) and $E_{(1,1)}E_{(2,2)}$.
\item $J^{(1)}$ and $J^{(2)}$ commute with $\lint^{(1)}$, $\lint^{(2)}$, the eight graded and the eight ungraded interior gammas, $H_g$, $H_u$ and the
interior copy rotation $M_{12}=\half\sum_\mu\eta_{\mu\mu}E_{(1,\mu)}E_{(2,\mu)}$. All of these are $\C$-linear on $W$ and on $\overline W$.
\item On $W$ the surplus part of the copy rotation vanishes, since $E_{(1,3)}E_{(2,3)}=E_{(1,2)}E_{(2,2)}\omega_S$, so the copy rotation
$C_b=\half\sum_a\eta_{aa}E_{(1,a)}E_{(2,a)}$ of Proposition~7.3(ii) of \cite{RE} acts on $W$ as $M_{12}$, $\C$-linearly. The four cross surplus bivectors $E_{(1,s)}E_{(2,t)}$
($s,t\in\{2,3\}$) preserve $W$ but anticommute with $J^{(1)}$ and $J^{(2)}$, so they act on $W$ antilinearly.
\end{enumerate}
\end{proposition}

\begin{proof}
(i) Bivectors of distinct constituents commute; $(J^{(1)}J^{(2)})^2=(J^{(1)})^2(J^{(2)})^2=\id$ and $\Tr(J\otimes J)=(\Tr J)^2=0$. $J^{(1)}\psi=J^{(2)}\psi$ is
equivalent to $J^{(1)}J^{(2)}\psi=-\psi$. $J^{(1)}-J^{(2)}$ is antisymmetric, so its kernel is the orthogonal complement of its image, and
$\omega_S(J^{(1)}-J^{(2)})=J^{(1)}-J^{(2)}$. The canonical map $S\otimes_\R S\to S\otimes_\C S$ has kernel $\operatorname{im}(J^{(1)}-J^{(2)})$, spanned by
$Jx\otimes y-x\otimes Jy$, so it restricts to an isomorphism on $W$; the case of $\overline W$ is the same with $-J$ on the second factor. (ii)
By equation~(3) of \cite{ES17} a blade commutes with the surplus four-blade $\omega_S$ iff it shares an even number of indices with it, and
otherwise anticommutes. (iii) The blade rule again: none of these contains a surplus index, except $J^{(k)}$ itself and $M_{12}$, whose
terms share no index with $J^{(1)}$ or $J^{(2)}$. (iv) $E_{(1,2)}E_{(2,2)}\omega_S=-E_{(2,2)}E_{(1,3)}E_{(2,2)}E_{(2,3)}=E_{(1,3)}E_{(2,3)}$, by the anticommutation of distinct
generators. On $W$, $\omega_S=-\id$, so the two surplus terms of $C_b$ cancel. A cross surplus bivector shares one index
with each $J^{(k)}$.
\end{proof}

The tensor-product linear algebra of (i) is inherited. Parts (ii)--(iv) are fixed (e1-3a, e1-3b, e1-3f). With the external $i$ of
Proposition~8.11(i) of \cite{RE}, $J^{(1)}+J^{(2)}$ is $2J$ on $W$ and $0$ on $\overline W$. So $W$ carries total charge $\pm2$, two particles
of equal charge, and $\overline W$ the neutral pairs, among them the particle--antiparticle state $\psi\otimes C\phi$ of that proposition
(restated; the charge reading interpretive). The question is now whether anything in the algebra selects $W$ and the complex
structure on it.

\begin{proposition}[Which imaginary unit the second rung selects]
\label{e1:prop:select}
\begin{enumerate}[label=(\roman*),leftmargin=2.2em]
\item Per-particle selection. The commutant of the eight ungraded gammas $\gamma^\mu\otimes\id$, $\id\otimes\gamma^\mu$ is $\HH\otimes\HH\cong M_4(\R)$,
spanned by $16$ blades. Adding the two temporal volume elements $T^{(1)},T^{(2)}$ (the datum of Remark~4.2 of \cite{ES17} for each copy)
cuts it to $\operatorname{span}\{\id,J^{(1)},J^{(2)},\omega_S\}\cong\C\otimes_\R\C\cong\C\oplus\C$. Its solutions of $x^2=-\id$ are exactly
$\pm J^{(1)}$ and $\pm J^{(2)}$, four isolated points, none of which acts as $J$ on both factors of $S_2$.
\item The spectator action, $k_I=1$. The selected unit of Theorem~4.1 of \cite{ES17} at $n=2$, the product of the eight internal generators, is
$J_{\rm sel}=J\otimes\omega=J^{(1)}\Gamma^{(2)}$ (Proposition~7.11(iv) of \cite{RE}). It commutes with $J^{(1)}$ and $J^{(2)}$ and so preserves
$W$. It equals $J^{(1)}$ on $\{\Gamma^{(2)}=1\}$ and $-J^{(1)}$ on $\{\Gamma^{(2)}=-1\}$, and $J^{(2)}$ on $\{\Gamma^{(2)}\omega_S=-1\}$ and
$-J^{(2)}$ on $\{\Gamma^{(2)}\omega_S=1\}$, $32$ dimensions each; on $W$ it is $J$ on $16$ and $-J$ on $16$ dimensions. It anticommutes with every generator of copy~2, graded or ungraded. So no gamma matrix of copy~2 is $J_{\rm sel}$-linear,
and neither $H_g$ nor $H_u$ is for $m_2\neq0$; $m_1E_{(1,1)}$ is.
\item The diagonal action, $k_I=2$. The commutant of the eight graded interior gammas $E_{(k,\mu)}$ is spanned by $16$ blades ($M_2(\HH)$). Adding
the six surplus bivectors cuts it to $\operatorname{span}\{\id,\omega_S\}\cong\R\oplus\R$, with $\omega_S^2=+\id$. No complex structure commutes with
the gammas of both copies and the surplus bivectors. On $W$ alone the restrictions of the same operators have commutant $\R$.
\item Up the ladder, under the diagonal action of rung $n$ the same commutant is $\operatorname{span}\{\id,\omega_S\}$ with
$\omega_S=J^{(1)}\cdots J^{(n)}$ and $\omega_S^2=(-1)^n$ (checked at $n=2,3$ for the commutant and $n=2,3,4$ for the square). An imaginary unit
is selected exactly at odd $n$. The common subspace $J^{(1)}=\cdots=J^{(n)}$, the complex $n$-fold product $S\otimes_\C\cdots\otimes_\C S$, has
dimension $2\cdot4^n$.
\end{enumerate}
\end{proposition}

\begin{proof}
(i) By the blade argument of \cref{e1:prop:copies}(ii), since $\id\otimes\gamma^\mu=\Gamma^{(1)}E_{(2,\mu)}$ and
$T^{(2)}=\Gamma^{(1)}E_{(2,1)}E_{(2,2)}E_{(2,3)}$ are blades. Equivalently, the commutant of the gammas on $S$ is the quaternion algebra of Remark~4.2
of \cite{ES17}, and that of $\{\gamma^\mu,T\}$ is $\R\oplus\R J$. For $x=a+bJ^{(1)}+cJ^{(2)}+d\,\omega_S$,
$x^2=(a^2-b^2-c^2+d^2)+2(ab-cd)J^{(1)}+2(ac-bd)J^{(2)}+2(ad+bc)\omega_S$, and $x^2=-\id$ has exactly the four solutions listed. (ii) $J_{\rm sel}$ is a blade of
grade $8$ containing every copy-2 index, so it anticommutes with each $E_{(2,a)}$; it commutes with $\Gamma^{(1)}$, so it also
anticommutes with each $u_a=\Gamma^{(1)}E_{(2,a)}$. $\Gamma^{(2)}$ commutes with
$J^{(1)}$ and $J^{(2)}$, and $\Tr\Gamma^{(2)}=0$, also on $W$. Finally $J_{\rm sel}\mp J^{(2)}=J^{(1)}(\Gamma^{(2)}\pm\omega_S)$ with $J^{(1)}$
invertible, and $\Gamma^{(2)}\omega_S$ is a traceless involution. (iii) A blade commuting with all eight interior generators meets them in all or none
of them. Meeting all of them requires odd grade, and meeting none requires even grade. Commuting with every surplus bivector forces it to meet
the four surplus directions in all or none. Both conditions together leave $\id$ and $\omega_S$. The algebra generated contains the
central idempotent $\half(\id-\omega_S)$, so its commutant on $W$ is the compression of $\operatorname{span}\{\id,\omega_S\}$, which is $\R$.
(iv) The same parity count with $4n$ interior and $2n$ surplus generators: a blade containing all $4n$ interior generators has grade $4n$ or
$6n$, both even, so it fails. The square follows from equation~(4) of \cite{ES17}, and the dimension is $8^n/2^{n-1}$, because the $n-1$
involutions $J^{(k)}J^{(k+1)}$ commute and have traceless products.
\end{proof}

Part (i) is fixed, with the quaternion commutant of one copy restated from Remark~4.2 of \cite{ES17}. Part (ii) restates Theorem~4.1 of
\cite{ES17} and Proposition~7.11(iv) of \cite{RE} (fixed there), and the comparison with $W$ is fixed. Part (iii) is obstructed under the
stated hypotheses (e1-3c to e1-3e, e1-3g). Part (iv) restates the selection at odd $n$ of Proposition~10.10(ii) of \cite{ES17} (fixed there)
and adds that nothing is selected at even $n$ (fixed; e1-3h).

\begin{corollary}[No selected two-particle complex structure]
\label{e1:cor:nocx}
At the second rung none of the selection principles of \cite{ES17} gives a complex structure under which both copies are complex Dirac spinors
with $\C$-linear gamma matrices. Per copy they give $\pm J^{(1)}$ and $\pm J^{(2)}$ separately. The spectator action gives
$J_{\rm sel}$, under which copy~2 carries no linear Dirac operator, so that copy~2 is internal. The diagonal action gives none. What the
diagonal action does select is the splitting $S_2=W\oplus\overline W$, the eigenspaces of the one nontrivial element $\omega_S$ of its commutant.
The complex structure $J|_W$ that makes $W$ the two-particle space commutes with every gamma matrix, with both Hamiltonians and with the copy
rotation, but not with the surplus $\so4$. It is an added datum.
\end{corollary}

The corollary is fixed as a statement about these selection principles (\cref{e1:prop:select}). That the choice of $W$, rather than
$\overline W$, and of $J|_W$ is what turns the second rung into a two-particle space is interpretive. This also bears on the phrase
``two particles from one''. The paper's own second quantisation of one constituent (Theorem~5.7 of \cite{ES17}, Proposition~3.9 of
\cite{RE}) already has two-excitation states of one Dirac mode, a level of dimension $6$ with charges $q=\pm2$ once each and $q=0$ four times.
In ordinary physics that is where several particles live (inherited \cite{Fock32}). It does not need the second rung.

\subsection{The Fock reading: quanta, not particles}
\label{e1:sec:fock}

The tower has its own number operator. $S_n$ is the Fock space of $3n$ fermionic modes, and $N$ counts them (Proposition~7.2 of
\cite{RE}). We ask which elements of $\so{3n,3n}$, the bivectors, conserve $N$.

\begin{proposition}[Bivectors and the number of quanta]
\label{e1:prop:Nbiv}
\begin{enumerate}[label=(\roman*),leftmargin=2.2em]
\item $N=\tfrac{3n}2-\half\sum_{k,i}E_{(k,i)}E_{(k,i+3)}$ is diagonal, with spectrum $0,\dots,3n$; at $n=2$ the multiplicities are
$1,6,15,20,15,6,1$. $N=N^{(1)}+N^{(2)}$, the same operator in the graded and ungraded readings. $(-1)^N=\Gamma^{(1)}\Gamma^{(2)}$ is the volume element of
$\Cl(6,6)$.
\item On vectors $[N,v]=K_nv$. A bivector $X$ commutes with $N$ iff $\ad X$ commutes with $K_n$. These form
$\mathfrak{gl}(L_+)\cong\mathfrak{gl}(3n,\R)$, of dimension $9n^2$ ($36$ at $n=2$), the stabiliser of the transverse maximal isotropic subspaces
$L_\pm=\ker(K_n\mp\id)$. The remaining $3n(3n-1)$ directions ($30$ at $n=2$) are $\Lambda^2L_+\oplus\Lambda^2L_-$, on which $[N,X]=\pm2X$: they create
or annihilate a pair of quanta.
\item $(-1)^N$ commutes with every bivector and anticommutes with every generator. No element of $\so{3n,3n}$ changes $N$ by an odd number.
\end{enumerate}
\end{proposition}

\begin{proof}
(i) $N_i=\half(\id-e_ie_{i+3})$, and the pole operators are diagonal (Appendix~A of \cite{ES17}). (ii) $[N,a^\pm_{k,i}]=\pm a^\pm_{k,i}$, and
$a^\pm_{k,i}$ is proportional to $e_i\pm e_{i+3}$ in constituent $k$, the $\pm1$ eigenvectors of $K$. $N-\tfrac{3n}2$ is a bivector, so $[N,X]$
is one, and the spin representation of $\so{V_n}$ is faithful. Hence $[N,X]=0$ iff $[\ad N,\ad X]=[K_n,\ad X]=0$. Under $\ad K_n$,
$\so{V_n}\cong\Lambda^2V_n=\Lambda^2L_+\oplus(L_+\otimes L_-)\oplus\Lambda^2L_-$ has eigenvalues $2,0,-2$, and $L_\pm$ are isotropic because
$\eta(e_i\pm e_{i+3},e_j\pm e_{j+3})=\delta_{ij}-\delta_{ij}$. (iii) Grade parity.
\end{proof}

Part (i) restates Proposition~7.2 of \cite{RE}. The rest is the standard structure of the spin representation on the exterior algebra of
a maximal isotropic subspace \cite{Chevalley54}, inherited, with its realisation in the tower fixed (e1-4a, e1-4b, e1-4g). The question
is which of these directions are physical.

\begin{theorem}[Compact fields conserve the quanta]
\label{e1:thm:Nrule}
Let $X\in\Cent(G_n)\cap\so{3n,3n}$, the algebra $\so{3n,\C}$ of binding-compatible bivectors (Theorem~17.19 of \cite{ES17}; Proposition~7.8 of
\cite{RE}). Then $K_nXK_n=-X^{\!\top}$, the Cartan involution. Consequently:
\begin{enumerate}[label=(\roman*),leftmargin=2.2em]
\item $X$ commutes with $N$ iff $X$ is compact (antisymmetric), that is iff it is a paired rotation $\bigl(\begin{smallmatrix}R&0\\0&R\end{smallmatrix}\bigr)$.
These form $\so{3n}$, of dimension $\tfrac32n(3n-1)$.
\item The non-compact part, the paired boosts $\bigl(\begin{smallmatrix}0&S\\-S&0\end{smallmatrix}\bigr)$, anticommutes with $K_n$. A nonzero paired
boost has nonzero components in both $\Lambda^2L_+$ and $\Lambda^2L_-$, so it changes $N$ by $+2$ and by $-2$ and conserves no part of it.
\item By Corollary~2.4 of \cite{ES17} no gauge boson is a boost. Hence no gauge-admissible, binding-compatible field creates or annihilates a
pair of quanta, at any rung.
\end{enumerate}
\end{theorem}

\begin{proof}
On vectors $\eta G_n=-K_n$ and $G_n\eta=-\eta G_n$. For $X$ commuting with $G_n$,
$K_nXK_n=\eta G_nX\eta G_n=-\eta G_nXG_n\eta=-\eta XG_n^2\eta=\eta X\eta$, and $\eta X\eta=-X^{\!\top}$ for $X\in\so{\eta}$. (i) follows from \cref{e1:prop:Nbiv}(ii). The block
forms are those of Proposition~7.8(i) of \cite{RE}. (ii) $X$ maps $L_\pm$ to $L_\mp$, so its components lie in $\Lambda^2L_+\oplus\Lambda^2L_-$.
$G_n$ commutes with $X$, maps $L_+$ onto $L_-$ (because $G_nK_n=-K_nG_n$) and normalises $\so{\eta}$, so conjugation by $G_n$ carries the $\Lambda^2L_+$
component of $X$ to its $\Lambda^2L_-$ component. They vanish together. (iii) Gauge admissibility is compactness, the criterion of
Section~2.2 of \cite{ES17}.
\end{proof}

The theorem is fixed. It was checked exactly on vectors at $n=1,\dots,4$, where $\dim\Cent(K_n)\cap\so{3n,3n}=9,36,81,144$,
$\dim\Cent(G_n)\cap\so{3n,3n}=6,30,72,132$ and the compact part has half of it. It was checked entry by entry on $S_n$ at $n=1,\dots,4$ for every
paired rotation $e_ie_j-e_{i+3}e_{j+3}$ and every paired boost $e_ie_{j+3}-e_je_{i+3}$, within one constituent and across two (e1-4c,
e1-4d). At $n=4$, $S_4$ has dimension $4096$, and the generators were handled as signed permutations.

\begin{proposition}[The quanta and the Lorentz algebras]
\label{e1:prop:Nlorentz}
The number $N$ behaves as follows under the copy algebra and under the two Lorentz algebras of one constituent.
\begin{enumerate}[label=(\roman*),leftmargin=2.2em]
\item The copy rotation $C_b$ commutes with the total number $N$, but it does not commute with the number $N^{(1)}$ of the first
constituent. The quarter turn
\[
U=e^{\pi C_b/2}=\tfrac18\textstyle\prod_a(\id+\eta_{aa}E_{(1,a)}E_{(2,a)})
\]
satisfies $UN^{(1)}U^{-1}=N^{(2)}$. The spin lift of the copy boost $\varepsilon\otimes G$ changes $N$ by $\pm2$ only, with both signs.
\item In $\lint$ the elements commuting with $N$ form $\R\,e_1e_4$, a boost (the pole operator $\id-2N_1$). The rotations $e_4e_5,e_4e_6,e_5e_6$, the
boosts $e_1e_5,e_1e_6$ and the charge $J=e_2e_3$ do not commute with $N$.
\item In $\lG$ the elements commuting with $N$ form its compact $\so3$, the paired rotations. No nonzero boost of $\lG$ commutes with $N$.
\end{enumerate}
\end{proposition}

\begin{proof}
(i) $C_b$ is a paired rotation (\cref{e1:thm:Nrule}(i)). Its six terms commute and square to $-\tfrac14$, so
$e^{\theta C_b}=\prod_a(\cos\tfrac\theta2+\sin\tfrac\theta2\,\eta_{aa}E_{(1,a)}E_{(2,a)})$, which at $\theta=\tfrac\pi2$ is $U$; it carries $E_{(1,a)}$ to
$\pm E_{(2,a)}$ (Proposition~7.3(ii) of \cite{RE}). The copy
boost is a paired boost (Proposition~7.8(iii) of \cite{RE}); apply \cref{e1:thm:Nrule}(ii). (ii) $\lint\cap\mathfrak{gl}(3,\R)$ is computed: of
the six planes in $\langle e_1,e_4,e_5,e_6\rangle$ only $\langle e_1,e_4\rangle$ is $K$-invariant, and $J$ involves the modes $a^+_2a^+_3$. (iii) is
\cref{e1:thm:Nrule} at $n=1$, with $\lG=\Cent(G)\cap\so{3,3}$ (Theorem~12.1 of \cite{ES17}).
\end{proof}

All three parts are fixed (e1-4e, e1-4f), and (i) gives a precise sense to the copy rotation as a hop of quanta between the constituents
at fixed total number. Part (ii) corrects an easy expectation. The Fock number is not invariant under the spatial rotations of the
interior four-plane, only under one of its boosts, and it does not commute with the charge $J$. The quanta counted by $N$ are therefore
not Lorentz-invariant particles. They are invariant under the compact part of $\lG$, the rotations of the field-strength reading, and under
nothing that contains a boost or the electric charge. Reading $N$ as a particle number is an internal reading tied to the split structure
$K$ (interpretive). What \cref{e1:thm:Nrule} does establish is the following. If quanta are counted by $N$, the only processes that create
them in pairs from the binding-compatible algebra are the paired boosts. These are excluded as gauge fields, so pair creation of quanta
would need a non-compact, non-gauge interaction (fixed, with the identification of quanta interpretive). This agrees with the
statement of Section~8.5 of \cite{RE} that duplication in the tower's sense is not pair creation.

\subsection{Exchange and statistics}
\label{e1:sec:exchange}

Two swaps act on $S_2$: the ungraded swap $P(x\otimes y)=y\otimes x$, and the Koszul swap $\tau_K=D\,P$ of Proposition~7.3(iii) of \cite{RE}, where
$D=(-1)^{\mathsf p\otimes\mathsf p}$ and $\omega=(-1)^{\mathsf p}$ on the basis of equation~(1) of \cite{ES17}.

\begin{proposition}[Two exchanges]
\label{e1:prop:swap}
\begin{enumerate}[label=(\roman*),leftmargin=2.2em]
\item $P(e_a\otimes\id)P=u_a=\Gamma^{(1)}E_{(2,a)}$. So $P$ exchanges the ungraded copies and maps the vectors $E_{(1,a)}$ of $\Cl(6,6)$ to elements of
grade $7$. It commutes with the diagonal Lorentz generators $\Delta_{\mu\nu}$, with $N$, with $\omega_S$ (so it preserves $W$) and with $H_u$ at
$m_1=m_2$, but not with $H_g$.
\item $\tau_KE_{(1,a)}\tau_K=E_{(2,a)}$. $\tau_K$ commutes with $\Delta_{\mu\nu}$, $N$, $\omega_S$ and $H_g$ at $m_1=m_2$, but not with $H_u$.
\item $D=\half(\id+\Gamma^{(1)}+\Gamma^{(2)}-\Gamma^{(1)}\Gamma^{(2)})$. So $\tau_K=P$ except on the sector $\Gamma^{(1)}=\Gamma^{(2)}=-1$, where $\tau_K=-P$.
\end{enumerate}
\end{proposition}

\begin{proof}
(i) $P(A\otimes B)P=B\otimes A$. Diagonal bivectors, $N=N^{(1)}+N^{(2)}$, $\omega_S=J\otimes J$ and $H_u$ at equal masses are symmetric under it,
while $P(\omega\otimes e_1)P=e_1\otimes\omega$. (ii) The first identity restates Proposition~7.3(iii) of \cite{RE}. Hence
$\tau_K\Gamma^{(1)}\tau_K=\Gamma^{(2)}$ and $\tau_Ku_1\tau_K=e_1\otimes\omega\neq u_1$, while bivectors $E_{(1,a)}E_{(1,b)}$ go to $E_{(2,a)}E_{(2,b)}$. (iii) $D$ is
diagonal with entry $-1$ exactly where both $\omega$-eigenvalues are $-1$.
\end{proof}

Each swap is a symmetry of its own reading and not of the other: $P$ of the additive Hamiltonian, $\tau_K$ of the graded one (fixed; e1-5a,
e1-5b). For the Lorentz content let $C=-\sum_{\mu<\nu}\eta^{\mu\mu}\eta^{\nu\nu}\Delta_{\mu\nu}^2$ be the diagonal Casimir, normalised as in Section~10.2
of \cite{ES17}. It is $\tfrac32$ on one constituent, it satisfies $C(C-3)(C-4)=0$ on $S_2$, and it takes the values $0$, $3$ and $4$ on
$(0,0)$, $(\half,\half)$ and $(1,0)\oplus(0,1)$. On one constituent each boost is $J\Gamma$ times its dual rotation,
$e_1e_4=J\Gamma\,e_5e_6$ and cyclically, so $\Gamma=\pm1$ are the two Weyl halves (e1-5c).

\begin{proposition}[Statistics sectors]
\label{e1:prop:stat}
The dimensions of the eigenspaces of $C$ in the sectors of $P$ and $\tau_K$ are those of \cref{e1:tab:stat}. In particular:
\begin{enumerate}[label=(\roman*),leftmargin=2.2em]
\item On $W$, $P$ is $\C$-linear. Its $-1$ eigenspace, $\Lambda^2_\C S$, carries $2(0,0)\oplus(\half,\half)$, and its $+1$ eigenspace, $\mathrm{Sym}^2_\C S$,
carries $(1,0)\oplus(0,1)\oplus(\half,\half)$. In the chirality sectors, $\Gamma^{(1)}=\Gamma^{(2)}=\pm1$ gives spin one in $P=+1$ and spin zero in
$P=-1$, while the mixed sector ($16$ dimensions, $C=3$) splits $8+8$.
\item On $W$, $\tau_K$ is $\C$-linear, and its two eigenspaces have the same Casimir content, $C=0,3,4$ on $2,8,6$ dimensions. The spin-one part of
$\tau_K=+1$ lies in $\Gamma^{(1)}=\Gamma^{(2)}=+1$, and that of $\tau_K=-1$ in $\Gamma^{(1)}=\Gamma^{(2)}=-1$: the two sectors are mirror images.
\item On $\overline W$, $P$ is antilinear, $PJ^{(1)}=-J^{(1)}P$, and $J^{(1)}$ exchanges its two eigenspaces. They carry the same content.
\end{enumerate}
\end{proposition}

\begin{proof}
All entries are traces of products of commuting integer matrices: the projectors $\half(\id\pm X)$ for $X\in\{\omega_S,P,\tau_K,\Gamma^{(1)},\Gamma^{(2)}\}$ and
the spectral projectors of $C$ (e1-5d to e1-5f). (i) $PJ^{(1)}P=J^{(2)}$, which equals $J^{(1)}$ on $W$. (ii) Use \cref{e1:prop:swap}(iii). (iii)
$J^{(2)}=-J^{(1)}$ on $\overline W$, and $J^{(1)}$ commutes with $C$.
\end{proof}

\begin{table}[t]
\centering
\small
\begin{tabular}{@{}lccc@{}}
\toprule
sector & on $W=S\otimes_\C S$ ($32$) & on $\overline W=S\otimes_\C\overline S$ ($32$) & on $S_2$ ($64$)\\
\midrule
$P=-1$ & $0^4\ 3^8$ & $0^2\ 3^8\ 4^6$ & $0^6\ 3^{16}\ 4^6$\\
$P=+1$ & $3^8\ 4^{12}$ & $0^2\ 3^8\ 4^6$ & $0^2\ 3^{16}\ 4^{18}$\\
$\tau_K=-1$ & $0^2\ 3^8\ 4^6$ & $0^2\ 3^8\ 4^6$ & $0^4\ 3^{16}\ 4^{12}$\\
$\tau_K=+1$ & $0^2\ 3^8\ 4^6$ & $0^2\ 3^8\ 4^6$ & $0^4\ 3^{16}\ 4^{12}$\\
\bottomrule
\end{tabular}
\caption{Casimir content $c^d$ ($C=c$ on $d$ real dimensions) of the exchange sectors under the diagonal interior Lorentz algebra.
Checks e1-5d to e1-5f.}
\label{e1:tab:stat}
\end{table}

The decompositions of $\Lambda^2$ and $\mathrm{Sym}^2$ of a Dirac spinor are inherited \cite{WeinbergQFT}; their appearance on $W$, and only on $W$,
is fixed. The Koszul column restates Proposition~7.10(iii) of \cite{RE} on $S_2$ (fixed there) and refines it on $W$. The two $\tau_K$
sectors are isomorphic as real representations, since the realification of $(1,0)$ is that of $(0,1)$, but they are mirror images once $J$ is
the imaginary unit (fixed). So the ungraded swap sorts the complex two-particle space by statistics in the textbook way. Identical fermions,
antisymmetric under $P$, have only scalar and vector content at a single momentum mode (inherited \cite{Pauli40}). The Clifford exchange
$\tau_K$ does not sort by Lorentz content at all. Its eigenvalue is a chirality-dependent sign, not a statistics label. Fermi statistics
for two copies is therefore a constructed requirement: antisymmetry under $P$ on $W$. The algebra supplies $\tau_K$, the exchange that
preserves the vectors of $\Cl(6,6)$, and $\tau_K$ is blind to it.

\subsection{Answer: what makes the second rung carry two particles}
\label{e1:sec:answer}

\Cref{e1:tab:answer} collects the structures. Read together they answer the question in four sentences. (1) As given by the tower, the second
rung carries one particle. Its vectors $E_{(k,a)}$ give one Dirac clock along a unit vector of $\R^{6,6}$, of mass $M=(m_1^2+m_2^2)^{1/2}$, with
no joint mass shell of the copies. Under the selected unit $J_{\rm sel}$ copy~2 carries no linear Dirac operator and is internal, so the rung
is one particle with a larger internal space (fixed; \cref{e1:prop:rest,e1:prop:select}). (2) The algebra also contains, uniquely up to
conjugation, a second commuting copy, the generators dressed by the chirality of the first. With it the copies have separate mass shells and
additive energies (fixed; \cref{e1:prop:copies,e1:prop:rest}). (3) Three further structures make that a two-particle system in the sense of
ordinary quantum mechanics: the subspace $W$ with the complex structure $J|_W$, the additive Hamiltonian $H_u$, and the ungraded exchange $P$
with antisymmetry. None of them is selected. $W$ is singled out only together with $\overline W$, and the diagonal action selects no imaginary unit at
even rungs (constructed; obstructed as a selection; \cref{e1:cor:nocx,e1:prop:stat}). (4) The tower's own counting operator $N$ counts quanta
of the split structure, not Lorentz-invariant particles. No gauge-admissible, binding-compatible field changes it (fixed;
\cref{e1:thm:Nrule,e1:prop:Nlorentz}).

\begin{table}[t]
\centering
\footnotesize
\begin{tabular}{@{}>{\raggedright\arraybackslash}p{3.6cm}>{\raggedright\arraybackslash}p{6.4cm}>{\raggedright\arraybackslash}p{4.6cm}@{}}
\toprule
structure & what it does at $n=2$ & class\\
\midrule
graded vectors $E_{(2,a)}$ & one Dirac clock of mass $M$; no joint mass shell; $\tau_K$ as exchange & fixed (they define $\Cl(6,6)$)\\
commuting copy $u_a=\Gamma^{(1)}E_{(2,a)}$ & two mass shells, energies $\pm m_1\pm m_2$ & fixed as an algebra (unique up to $\id\otimes GL(8,\R)$); use as vectors constructed\\
common complex structure $J|_W$ & $W=S\otimes_\C S$, $32$ dimensions, total charge $\pm2$ & $W\oplus\overline W$ fixed by $\omega_S$; $J|_W$ constructed; selection obstructed ($k_I=2$), copy~2 internal ($k_I=1$)\\
additive Hamiltonian $H_u$ & threshold $m_1+m_2$; two-torus rest flow & constructed; not a vector (grades $1$ and $7$)\\
ungraded exchange $P$ & $\Lambda^2_\C S=2(0,0)\oplus(\half,\half)$ on $W$ & constructed; $\tau_K$ fixed but statistics-blind\\
Fock number $N$ & counts quanta; conserved by paired rotations, changed by $\pm2$ by paired boosts & fixed; as a particle number interpretive (not Lorentz- or $J$-invariant)\\
\bottomrule
\end{tabular}
\caption{Structures that bear on a two-particle reading of the second rung.}
\label{e1:tab:answer}
\end{table}

\paragraph{Which process takes one to two.} Three candidates have been tested. The step from rung $1$ to rung $2$ is a change of arena,
not a dynamical process. In the graded reading it enlarges the particle and does not add one. In the Fock reading the Lie algebra
$\so{3n,3n}$ changes $N$ only by $0$ or $\pm2$, and among binding-compatible fields only the paired boosts do it, which are not gauge fields.
A step from one quantum to two would need an odd element, which is a field operator and not a symmetry (fixed;
\cref{e1:prop:Nbiv,e1:thm:Nrule}). In the two-particle reading the step is the tensor embedding $\psi\mapsto\psi\otimes\phi_0$, which intertwines
$A\mapsto A\otimes\id$ but needs a state $\phi_0$ of the second particle that nothing selects (constructed). In ordinary physics
the process is creation from a field, and the paper's second quantisation of one constituent already contains it (Theorem~5.7 of
\cite{ES17}). Whether some mechanism of the arena selects $W$ and $J|_W$ dynamically is open.

\paragraph{Up the ladder.} What has been proved at every rung is this. The commuting copies exist, dressed by $\Gamma^{(1)}\cdots\Gamma^{(k-1)}$.
The graded masses add in quadrature and the ungraded ones additively. The rule of \cref{e1:thm:Nrule} holds. The common complex subspace
$J^{(1)}=\cdots=J^{(n)}$ has dimension $2\cdot4^n$. The diagonal action selects an imaginary unit, $\omega_S=J^{(1)}\cdots J^{(n)}$, exactly at
odd $n$. On the common subspace it acts as $(-1)^{(n-1)/2}J$, so at odd rungs the selected unit and the common complex structure agree up to
sign there, while at even rungs there is no selected unit to compare (fixed; e1-1d, e1-2f, e1-3h, e1-4c, e1-4d). The physics of the odd rungs, whose diagonal sector is the vector-like matter of Section~10 of \cite{ES17},
is not examined here.

\begin{table}[t]
\centering
\footnotesize
\begin{tabular}{@{}>{\raggedright\arraybackslash}p{8.3cm}>{\raggedright\arraybackslash}p{4.4cm}>{\raggedright\arraybackslash}p{2.2cm}@{}}
\toprule
statement & class & checks\\
\midrule
\cref{e1:prop:copies}: commuting copy, commutant $\id\otimes M_8(\R)$, grade-$7$ generators, chirality dressing, rung $n$ & restated (RE Prop.~7.1, ES17 App.~A); dressing fixed; exponential law inherited & e1-1a to e1-1d\\
\cref{e1:prop:rest}: quadrature against addition, rotated clock, joint mass shells, rest flows, rung $n$ & fixed; restated at $m_1=m_2$ (RE Prop.~8.11(iv)); two-body kinematics inherited & e1-2a to e1-2g\\
\cref{e1:prop:W}: $W=S\otimes_\C S$, $\omega_S$, what preserves $W$, $C_b|_W=M_{12}|_W$ & fixed; tensor algebra inherited & e1-3a, e1-3b, e1-3f\\
\cref{e1:prop:select}: per-particle units $\pm J^{(k)}$; $J_{\rm sel}=J^{(1)}\Gamma^{(2)}$; no unit at $k_I=2$; parity of $n$ & fixed; restated (ES17 Thm~4.1, Prop.~10.10; RE Prop.~7.11); obstructed & e1-3c to e1-3e, e1-3g, e1-3h\\
\cref{e1:cor:nocx}: no selected two-particle complex structure & fixed (as stated); reading interpretive & ---\\
\cref{e1:prop:Nbiv}: $\ad N=K_n$, $\mathfrak{gl}(3n,\R)$, pairs, parity & restated (RE Prop.~7.2); inherited; realisation fixed & e1-4a, e1-4b, e1-4g\\
\cref{e1:thm:Nrule}: compact $\Leftrightarrow$ conserves $N$; boosts change it by $\pm2$; no gauge field creates pairs & fixed & e1-4c, e1-4d\\
\cref{e1:prop:Nlorentz}: copy rotation hops quanta; $N$ against $\lint$, $J$, $\lG$ & fixed; particle reading interpretive & e1-4e, e1-4f\\
\cref{e1:prop:swap}: $P$ against $\tau_K$ & fixed; $\tau_K$ restated (RE Prop.~7.3) & e1-5a, e1-5b\\
\cref{e1:prop:stat}, \cref{e1:tab:stat}: statistics sectors & fixed; Dirac$\otimes$Dirac inherited; $\tau_K$ on $S_2$ restated (RE Prop.~7.10) & e1-5c to e1-5f\\
\cref{e1:sec:answer}: which process; up the ladder; dynamical selection of $W$ & constructed; fixed; open & see text\\
\bottomrule
\end{tabular}
\caption{Claims of \cref{e1:sec:main} and their classes. The checks are the \texttt{PASS} lines printed by the script
\texttt{verify/e1\_two\_from\_one.py}.}
\label{e1:tab:claims}
\end{table}
